% ID: SEL-D0050-OB-E02
% ESTADO: validada. S051.
\textbf{a)} La parábola de $f'$ abre hacia arriba, pues su vértice está por debajo de sus dos raíces. Por ello, $f'>0$ en $(-\infty,-1)$ y $(5,\infty)$ y $f'<0$ en $(-1,5)$. La función $f$ crece en los dos primeros intervalos de signo positivo y decrece entre las raíces.

\textbf{b)} En $x=-1$ la derivada pasa de positiva a negativa, de modo que hay un máximo relativo. En $x=5$ pasa de negativa a positiva, así que hay un mínimo relativo.

\textbf{c)} El vértice da $f'(2)=-4$. Con el punto $(2,5)$, la ecuación es $y-5=-4(x-2)$, es decir, $\boxed{y=-4x+13}$.
